% CCSS 7.RP.A.2.A — Recognizing Proportional Relationships
\section{Recognizing Proportional Relationships}

\begin{learningGoals}
\begin{itemize}
    \item Decide whether two quantities are proportional using a table or graph
    \item Explain why a relationship is or is not proportional
\end{itemize}
\end{learningGoals}

\begin{conceptBox}{What Makes It Proportional?}
Two quantities are \textbf{proportional} if every pair has the \textbf{same ratio}.
\begin{itemize}[leftmargin=*]
    \item \textbf{Table test:} Divide $y \div x$ for every row. Same quotient every time $=$ proportional.
    \item \textbf{Graph test:} Points form a \textbf{straight line through the origin} $(0, 0)$.
\end{itemize}
One mismatch means not proportional.
\end{conceptBox}

\begin{workedExample}{Table Test}
\begin{center}
\begin{tabular}{|c|c|c|}
\hline
\textbf{Hours} ($x$) & \textbf{Pay} ($y$) & $y \div x$ \\
\hline
$2$ & $\$24$ & $12$ \\
$5$ & $\$60$ & $12$ \\
$8$ & $\$96$ & $12$ \\
\hline
\end{tabular}
\end{center}
Every ratio is $12$ → \textbf{proportional} at $\$12$/hr.
Now add a $\$20$ sign-up fee: costs become $\$44, \$80, \$116$. Ratios: $22, 16, 14.5$ → \textbf{not proportional}.
\end{workedExample}

\begin{errorBox}[Just a Straight Line?]
``The graph is a straight line, so it must be proportional.'' \textbf{Not quite!} The line must also pass through $(0, 0)$. A line crossing the $y$-axis at $5$ is not proportional.
\end{errorBox}

\mascotSays{No sign-up fees, no starting amounts. Proportional means it starts at zero and stays straight!}

\begin{practiceBox}{Recognizing Proportional Relationships}
\resetProblems

\prob Is this relationship proportional? Explain.
\begin{center}
\begin{tabular}{|c|c|}
\hline $x$ & $y$ \\ \hline $3$ & $9$ \\ $5$ & $15$ \\ $7$ & $21$ \\ \hline
\end{tabular}
\end{center}
\answerExplain{Yes}{Divide $y \div x$ for each row: $9 \div 3 = 3$, $15 \div 5 = 3$, $21 \div 7 = 3$. Every ratio is $3$, so the relationship is proportional.}

\prob Is this relationship proportional? Explain.
\begin{center}
\begin{tabular}{|c|c|}
\hline $x$ & $y$ \\ \hline $2$ & $6$ \\ $4$ & $10$ \\ $6$ & $14$ \\ \hline
\end{tabular}
\end{center}
\answerExplain{No}{Divide $y \div x$ for each row: $6 \div 2 = 3$, $10 \div 4 = 2.5$, $14 \div 6 \approx 2.3$. The ratios are not all equal, so this is not proportional. It has a constant change of $+4$, but that does not mean a constant ratio.}

\trueOrFalse{A straight line through $(0, 0)$ always shows a proportional relationship.}
\answerTF{True}

\prob A phone plan costs $\$0.05$ per text with no monthly fee. A second plan costs $\$5$/month plus $\$0.02$ per text. Which is proportional?
\answerExplain{Plan 1}{Plan 1 has no base fee, so dividing total cost by texts always gives $\$0.05$. That constant ratio means it is proportional. Plan 2 has a $\$5$ monthly charge that does not depend on texts, so the cost-per-text ratio changes and it is not proportional.}

\end{practiceBox}
